\scanpage{249}
\mspage{248}
\mspage{6}\hfill n$^\circ$ 293

$\unsure{z|\mathfrak z}$ soit l'ensemble des éléments $z\in\unsure{r|\mathfrak r}$ tels que $zu\in\unsure{f|\mathfrak f}$.

Posons $r'=\unsure{o|\mathfrak o}[t]+ur$ ; on vérifie immédiatement que $r'$ est un sous-anneau de $r$. Si $y$ est un élément de $o[t]$, une condition nécessaire et suffisante pour que $yr'\subset o[t]$ est que $yur\subset o[t]$, donc que $yu\in f$, ou encore que $y\in z\cap o[t]=\unsure{y|\mathfrak y}$; $y$ est donc bien le conducteur de $r'$ par rapport à $o[t]$.

Supposons pour un moment que $\omega(t)$ ne soit pas transcendant par rapport à $\omega(o)$ ; il y a alors un polynôme $\bar H\ne0$ à coefficients dans $\omega(o)$ tel que $\bar H(\omega(t))=0$. On en déduit tout de suite qu'il y a un polynôme $H$ à coefficients dans $o$ dont le coefficient dominant $c$ n'est pas dans $z$ tel que $H(t)\in z$. On a alors $H(t)\in y$ ; il résulte alors du lemme 3, appliqué au couple $(o,r')$, qu'il y a une puissance de $c$ qui appartient à $y$, donc à $z$ ; or ceci est impossible puisque $z$ est premier et $c\notin z$.

\uline{Lemme 5} - La \uline{situation étant celle décrite dans l'énoncé du lemme} 3,- \uline{supposons de plus que} $o$ \uline{soit noetherien, Soient} $\unsure{q|\mathfrak q}$ \uline{un idéal premier maximal de} $r$ \uline{et} $\unsure{m|\mathfrak m}$ \uline{l'idéal premier} $q\cap o$ \uline{de} $o$, \uline{Si l'anneau local de} $q$ \uline{est distinct de celui de l'idéal} $m$ \uline{de} $o$, \uline{il y a un idéal premier de} $r$ \uline{contenu dans} $q$, \uline{distinct de} $q$ \uline{et contenant} $m$.

Supposons d'abord que $q$ contienne le conducteur $f$ de $r$ par rapport à $o[t]$ ; nous allons voir qu'on est alors dans le second cas. L'idéal $\quad$ contient un élément minimal $\unsure{z|r}$ de l'ensemble des idéaux premiers qui contiennent $f$ ; $z$ est donc un idéal premier associé à $f$. Soit $\omega$ l'homomorphisme canonique de $r$ sur $r/z$ ; $\omega(t)$ est transcendant sur $\omega(o)$ (lemme 3) et $r/z$ est entier sur $(\omega(o))\,[\omega(t)]$. Le noyau de la restriction
